% Part: first-order-logic
% Chapter: sequent-calculus
% Section: quantifier-rules

\documentclass[../../../include/open-logic-section]{subfiles}

\begin{document}

\olfileid{fol}{seq}{qrl}

\olsection{પરિમાણકના નિયમો}

\subsection{$\lforall$ માટેના નિયમો}

\begin{defish}
\Axiom$ !A(t), \Gamma \fCenter \Delta$
\RightLabel{\LeftR{\lforall}}
\UnaryInf$ \lforall[x][!A(x)],\Gamma \fCenter \Delta$
\DisplayProof
\hfill
\Axiom$ \Gamma \fCenter \Delta, !A(a) $
\RightLabel{\RightR{\lforall}}
\UnaryInf$ \Gamma \fCenter \Delta, \lforall[x][!A(x)]$
\DisplayProof
\end{defish}

\LeftR{\lforall}માં $t$ બંધ પદ (એટલે કે ચલ વિનાનું પદ) છે.
\RightR{\lforall}માં $a$ એવું !!a{constant} છે જે
\RightR{\lforall} નિયમના નીચલા સિક્વન્ટમાં ક્યાંય ન આવવું જોઈએ.
\RightR{\forall} અનુમાનમાં $a$ને \emph{આઇગનચલ} કહીએ
છીએ.\footnote{ઉપરના નિયમમાં $a$ !!a{constant} હોવા છતાં આપણે
``આઇગનચલ'' શબ્દ વાપરીએ છીએ. તેનાં ઐતિહાસિક કારણો છે.}

\subsection{$\lexists$ માટેના નિયમો}

\begin{defish}
\Axiom$ !A(a), \Gamma \fCenter \Delta $
\RightLabel{\LeftR{\lexists}}
\UnaryInf$ \lexists[x][!A(x)], \Gamma \fCenter \Delta$
\DisplayProof
\hfill
\Axiom$ \Gamma \fCenter \Delta, !A(t) $
\RightLabel{\RightR{\lexists}}
\UnaryInf$ \Gamma \fCenter \Delta, \lexists[x][!A(x)]$
\DisplayProof
\end{defish}

અહીં ફરીથી $t$ બંધ પદ છે, અને $a$ એવું !!a{constant} છે જે
\LeftR{\lexists} નિયમના નીચલા સિક્વન્ટમાં આવતું નથી. \LeftR{\lexists}
અનુમાનમાં $a$ને \emph{આઇગનચલ} કહીએ છીએ.

\RightR{\lforall} અથવા \LeftR{\lexists} અનુમાનના નીચલા સિક્વન્ટમાં
આઇગનચલ ન આવવું જોઈએ, એ શરતને \emph{આઇગનચલ-શરત} કહે છે.

\begin{explain}
યાદ કરો કે જ્યારે $!A$ એવું !!a{formula} હોય જેમાં !!{variable}~$x$
મુક્ત હોય, ત્યારે આપણે $!A(x)$ લખીને તે દર્શાવીએ છીએ. એ જ સંદર્ભમાં
$!A(t)$ એ~$\Subst{!A}{t}{x}$નું સંક્ષિપ્ત રૂપ છે. તેથી
\RightR{\lexists} નિયમને આ પ્રમાણે પણ લખી શકીએ:
\begin{prooftree}
  \Axiom$\Gamma \fCenter \Delta, \Subst{!A}{t}{x}$
  \RightLabel{\RightR{\lexists}}
  \UnaryInf$\Gamma \fCenter \Delta, \lexists[x][!A]$
\end{prooftree}
નોંધો કે $t$ એ~$!A$માં પહેલેથી આવી શકે છે; દાખલા તરીકે, $!A$ એ
$\Atom{\Obj P}{t,x}$ હોઈ શકે. તેથી $\Gamma \Sequent \Delta,
\Atom{\Obj P}{t,t}$ પરથી $\Gamma \Sequent \Delta,
\lexists[x][\Atom{\Obj P}{t,x}]$નું અનુમાન કરવું એ
\RightR{\lexists}નો યોગ્ય ઉપયોગ છે—આધારવિધાનમાં આવતાં $t$નાં એક કે વધુ,
પરંતુ અનિવાર્યપણે બધાં જ નહિ, સ્થાનોને બંધાયેલા !!{variable}~$x$થી
``બદલી'' શકાય છે. જોકે, \RightR{\lforall} અને~\LeftR{\lexists}માંની
આઇગનચલ-શરતો માગે છે કે !!{constant}~$a$ એ~$!A$માં ન આવે. તેથી
$\Gamma \Sequent \Delta, \Atom{\Obj P}{a,a}$ પરથી
$\Gamma \Sequent \Delta, \lforall[x][\Atom{\Obj P}{a,x}]$નું
અનુમાન કરવા માટે~$\RightR{\lforall}$નો ઉપયોગ યોગ્ય રીતે થઈ શકતો નથી.
\end{explain}

\begin{explain}
\RightR{\lexists} અને \LeftR{\lforall}માં પદ~$t$ પર કોઈ પ્રતિબંધ નથી.
બીજી બાજુ, \LeftR{\lexists} અને \RightR{\lforall} નિયમોમાં
આઇગનચલ-શરત માગે છે કે !!{constant}~$a$ ઉપરના સિક્વન્ટમાં $!A(a)$ની બહાર
ક્યાંય ન આવે. તંત્ર યથાર્થ રહે—એટલે કે માત્ર પ્રામાણ્ય સિક્વન્ટો જ
!!{derive}s કરે—તેની ખાતરી માટે આ જરૂરી છે. આ શરત વિના નીચેનું માન્ય હોત:
\begin{prooftree}
  \Axiom$!A(a) \fCenter !A(a)$
  \RightLabel{*\LeftR{\lexists}}
  \UnaryInf$\lexists[x][!A(x)] \fCenter !A(a)$
  \RightLabel{\RightR{\lforall}}
  \UnaryInf$\lexists[x][!A(x)] \fCenter \lforall[x][!A(x)]$
  \DisplayProof\bottomAlignProof
  \qquad
  \Axiom$!A(a) \fCenter !A(a)$
  \RightLabel{*\RightR{\lforall}}
  \UnaryInf$!A(a) \fCenter \lforall[x][!A(x)]$
  \RightLabel{\LeftR{\lexists}}
  \UnaryInf$\lexists[x][!A(x)] \fCenter \lforall[x][!A(x)]$
\end{prooftree}
પરંતુ $\lexists[x][!A(x)] \Sequent \lforall[x][!A(x)]$ પ્રામાણ્ય નથી.
\end{explain}

\end{document}
